\section{DISCUSSION}
Our investigation into how BVI individuals learn video creation through structured workshops reveals fundamental challenges in translating visual media production into accessible pedagogy.  By examining barriers and collaborative patterns, we identify how access is achieved not by universal solutions, but through dynamic, localized negotiations among individuals with different sensory abilities . We organize our conversation around three main themes: (1) what inclusive video creation education requires in terms of materials, pacing, and instructional approaches, (2) how collaboration emerges through dynamic negotiations between participants with different sensory abilities, and (3) how distributed expertise across peer networks can compensate for individual knowledge gaps.

\subsection{Implications for Inclusive Video Creation Education}

\subsubsection{Reframing Access Through Interaction Design }
Our findings demonstrate that the fundamental challenge in ability-diverse video creation education lies not in the epistemic differences between visual and non-visual ways of knowing, but in the interaction asymmetries these differences create. While participants with different vision profiles inevitably construct distinct mental models of video creation, for example, visual-spatial for sighted and low-vision users, sequential-temporal for screen reader users, these epistemic frameworks need not be barriers to learning. As our collaborative patterns reveal, participants successfully navigated these differences when interaction barriers were addressed through haptic mediation (Pattern 1) and cross-modal translation (Pattern 2).

This distinction between epistemic and interaction asymmetry extends Wobbrock et al.'s ability-based design framework \cite{wobbrock_ability-based_2011} by suggesting that educational technologies should not attempt to normalize different ways of knowing but rather create interaction bridges between them. For instance, when the low-vision tutor in Episode 2-1 served as a "modality relay," translating gestural instructions into VoiceOver-compatible actions, they preserved the integrity of both epistemic frameworks while enabling successful task completion. This suggests that future video creation tools designed for education should prioritize interaction flexibility over content uniformity, allowing multiple valid pathways to achieve creative goals.

\subsubsection{From Individual to Dyadic Pedagogies }
Perhaps our most significant finding concerns the reconceptualization of the learner unit in accessible media education. Rather than treating BVI participants as individual learners requiring accommodation, our data reveals the collaborative dyad as a more effective pedagogical unit. Episodes 1-1 and 1-2 demonstrate how helper-participant pairs developed sophisticated collaborative protocols that went beyond simple assistance to create what Bennett et al. term "interdependence \cite{bennett_interdependence_2018}."  However, our findings extend this framework by showing that interdependence can be explicitly taught and scaffolded rather than emerging spontaneously.

Educating collaborative dyads would involve several strategic shifts. First, orientation sessions should address both participants and their support workers together, establishing shared vocabulary and interaction protocols. Second, curricula should include explicit instruction in cross-modal communication strategies, teaching helpers how to provide spatial information through haptic guidance and verbal description that aligns with screen reader navigation. Third, assessment criteria should evaluate collaborative achievement rather than individual mastery, recognizing that distributed expertise represents a valid and valuable form of media literacy. This approach transforms support workers from peripheral accommodations into integral pedagogical partners, a shift that acknowledges the reality of collaborative creative practice in professional media production.

\subsubsection{Infrastructural Requirements for Ability-Diverse Learning }
The tool and ecosystem asymmetries we identified point toward specific infrastructural requirements for sustainable video creation education. Our participants encountered not app-level barriers but workflow-level friction when screen reader navigation and editing paradigms misaligned. This suggests that accessible video education requires standardized, cross-platform tools that maintain consistency across devices and operating systems. Web-based applications present a promising approach for addressing platform fragmentation, as they could provide uniform interfaces across iOS and Android while leveraging established web accessibility standards.

However, given the current absence of fully accessible video editing platforms, education programs must develop systematic workarounds. Our findings suggest these workarounds should be formally documented in multiple accessible formats—text descriptions for screen reader users, visual guides for magnifier users, and audio walkthroughs for review between sessions. The documentation gap P4 identified ("it became a memory exercise") could be addressed through persistent, multi-modal learning artifacts that participants can reference during independent practice.

\subsubsection{Cultivating Peer Expertise Networks }
The emergence of sophisticated peer teaching in our remote sessions (Pattern 4) reveals an underutilized resource in accessible media education. The "flatter" structure of online workshops, where participants occupied equal visual space in gallery view without the hierarchical arrangement of traditional classrooms, facilitated peer-to-peer knowledge exchange. Episodes 4-1 demonstrate how BVI participants with advanced technical skills provided precise spatial guidance and cross-version label mapping that preserved learner agency while ensuring task completion.

This peer expertise should be intentionally cultivated through structured opportunities for knowledge sharing. Mixed-experience cohorts, where novice and advanced BVI creators learn together, can create natural mentorship relationships. Formal peer teaching roles, where experienced participants demonstrate specific techniques or troubleshoot common problems, validate BVI expertise while providing relatable instruction. Remote learning environments particularly support this approach, as they eliminate physical barriers to peer assistance while maintaining equal participant presence. Additionally, implementing screen sharing and remote control capabilities in online sessions could address the limitation we observed where helpers could only provide verbal guidance, extending the haptic mediation strategies successful in face-to-face contexts to digital environments.

\subsection{From Universal Design to Dynamic Standard-Making}
Our results challenge the assumption that universal design can sufficiently accommodate ability-diverse learning environments. The three asymmetries we identified: expertise, tools, and interaction paradigms, demonstrate that participants function within fundamentally divergent epistemic frameworks based on their sensory modalities. This corresponds with theories of embodied cognition \cite{varela_embodied_1991,wilson_six_2002} illustrate that various sensory modalities engender unique cognitive frameworks.

The visual-touch, visual-assisted, and voice-tactile modes we observed represent not simply different ways of accessing the same information, but different ways of knowing and interacting with the world. This aligns with Thieme et al.'s research \cite{thieme_enabling_2017} on how individuals with visual impairments navigate their capabilities in social settings, and it expands upon Moser's assertion that disability is manifested differently across socio-material situations \cite{moser_disability_2006}.

This epistemic variety poses a challenge to the concepts of Universal Design for Learning (UDL) \cite{inc_udl_2024,rose_universal_2001}, which assume that various modes of representation, interaction, and expression may accommodate all learners. Although UDL has advanced inclusive education significantly \cite{courey_improved_2013,rao_using_2016}, our research indicates that its framework may unintentionally enforce a unique epistemic standard. Similar critiques have emerged from disability studies scholars who argue that universal design often defaults to normative assumptions \cite{dolmage_academic_2017,hamraie_building_2017}. 

Instead, our collaboration patterns illustrate what we refer to as "dynamic standard-making": the creation of temporary, localized standards between specific individuals that enable particular learning moments. This idea is based on Suchman's situated action theory \cite{suchman_human-machine_2006}, It says that standards come from practice, not the other way around, also Wobbrock et al.'s ability-based design framework \cite{wobbrock_ability-based_2011}, which argues that systems should adapt to users' abilities rather than requiring users to adapt to systems. Our findings extend this by showing that in collaborative contexts, the "system" itself is co-constructed through practice. When participants engaged in haptic mediation or cross-modal translation, they were not just adapting to each other's abilities but actively constructing what Lave and Wenger call "communities of practice" \cite{lave_situated_1991} with their own internal logics. These temporary standards represent ability-based design enacted socially, where the collaborative dyad or triad creates custom interaction protocols that leverage each person's specific abilities in that moment.

\subsection{The Distribution of Access}
Building on Goodwin's concept of "co-operative transformation zones," \cite{goodwin_co-operative_2013} we see how participants progressively transform the substrate of prior interactions to create new forms of collaboration. Each workshop session builds upon resources and techniques developed in previous sessions, creating what Goodwin calls an "accumulative" process central to human culture and knowledge.

The progression from direct haptic mediation to sophisticated peer teaching exemplifies what Due terms "distributed perception" \cite{due_distributed_2021}: where perception-related actions are shared among agents who co-operate to accomplish activities that could not be achieved individually. In our workshops, BLV participants relied on sighted participants' visual perception, but crucially, this was not passive dependency. Instead, as Due demonstrates, distributed perception requires active "positioning for perception" \cite{goodwin_participation_2007}, where participants must bodily and technologically position themselves to enable perception-related actions.

This distributed model challenges traditional special education approaches where expertise flows unidirectionally from teacher to student \cite{giangreco_students_2002}. Instead, we observed what Bennett et al. term "interdependence," \cite{bennett_interdependence_2018} where all parties contribute essential knowledge. The evolution of collaborative patterns we documented suggests that groups develop what Akkerman and Bakker call "boundary crossing competence," \cite{akkerman_boundary_2011} becoming increasingly skilled at navigating between different epistemic frameworks. 

This distributed model contrasts with research on paraprofessional support showing that traditional helpers often create dependency rather than independence \cite{causton-theoharis_increasing_2005}. Our observations align better with Bennett et al.'s interdependence framework \cite{bennett_interdependence_2018} and recent work on collective access \cite{mia_mingus_access_2017}, where access is understood as a collective responsibility rather than individual accommodation.
